\SourceLeaf{351}
In my (the first) construction, where $\delta$ is variable, the clise is the same as that of
\[
(a-b)(a-c)(b-c).
\]
If we denote by $e$ the fourth intersection [point] of the ellipse with the circle through $b$, $c$, $d$ and set
\[
e = \varepsilon m + \frac{n}{\varepsilon},
\]
which gives
\[
\beta\gamma\delta\varepsilon = 1
\]
then, in \textsc{Newton}'s construction, the clise of the difference of two values of $t$ is the same as that of
\[
(c-b)(b-e)(c-e).
\]
\Needspace{10\baselineskip}The above formula ($a=\varepsilon\rho$) can also be written better as follows:
\[
t = -\frac{(m-\delta\beta n)(m-\delta\gamma n)(m-\delta\varepsilon n)\rho}{\delta(1-\rho)mn} + \frac{(\beta+\gamma)mm-\dfrac{1}{\varepsilon}mn+\delta nn}{m}
\]
\[
= -\frac{(b-c)(b-e)(c-e)}{\delta(\beta-\gamma)(\beta-\varepsilon)(\gamma-\varepsilon)}\cdot\frac{\rho}{(1-\rho)mn} + \frac{(\beta+\gamma)mm-\dfrac{1}{\varepsilon}mn+\delta nn}{m}.
\]

\Needspace{10\baselineskip}\begin{center}\textbf{REMARKS.}\end{center}

Note [III], which is on a separate leaf, reproduces the substance of investigations that \textsc{Gauss} reported in his letter printed on pp.\,295 to 296 of this volume, addressed to \textsc{Schumacher} and dated 12 May 1843; the reader may refer to it for an explanation of the text. Notes [I] and [II], which are closely connected with [III], are found in a handbook and on a slip of paper; [I] probably dates as far back as 1831, [II] presumably from around 1840.
\hfill \textsc{Stäckel}.
